%document class
\documentclass[10pt,oneside]{book}
%%%% Page Info + Commands %%%%%{

%packages
\usepackage{pgfplots}
\pgfplotsset{compat=1.18}
\usepackage{amsmath}
\usepackage{amsfonts}
\usepackage{amsthm,letltxmacro,changepage}
\usepackage{enumerate}
\usepackage{enumitem}
\usepackage{changepage}

%This will shrink the page
\newcommand{\smallpage}{  \setlength \oddsidemargin{.5in}
            \setlength \textwidth{5in}
            \setlength \topmargin{0in}
            \setlength \textheight{9in}}


% Commands for sets of numbers
\newcommand{\Z}{\mathbb{Z}}\newcommand{\R}{\mathbb{R}}\newcommand{\N}{\mathbb{N}}

\newcommand{\n}{\vspace{0.15cm}\\}
\newcommand{\en}{\vspace{0.25cm}\\}

\newcommand{\gimage}[3]{
\begin{figure}[h]   
    \centering          
    \includegraphics[width=#2\textwidth]{#1}  
    \caption{#3}
    \label{fig:#1}
\end{figure}\\
}

\newcommand{\centerit}[1]{{\begin{center} #1 \end{center}}}
\newcommand{\ul}[1]{\underline{#1}}

%%%%%%%% command for graphics %%%%%%%%%%%%%
\usepackage{fancyhdr}
%}

%%%% Page Setup %%%%%%%
\smallpage\sloppy
\pagestyle{fancy}
\lhead{\large{\textbf{Day 1 - Intro to Proofs}}}
%\rfoot{\thepage/\pageref{LastPage} }
\setlength{\headheight}{10pt}


% Proof writing
\LetLtxMacro\oldproof\proof
\let\endoldproof\endproof
\renewenvironment{proof}[1][\proofname]
  {\vspace{0.2cm}\begin{adjustwidth}{2em}{2em}
   \oldproof[#1]}
  {\endoldproof
   \end{adjustwidth}}


\begin{document}

%%%% Introductions
\section{Introductions}
Class begins by discussion \textit{how we prove something}. A proof is basically a \textit{really really really small paper}, and the point of communicatino is to tell a story in a way. We should sit and think, ``how is this going to go?''
\subsection{Statements}
A statement is something that has a \textbf{truthness} or \textbf{falseness} to it. For example, 
\[\text{The sky is blue today}\]
It may be impossible to discern a statement's truth value, but that doesn't mean that it's not a statement (e.g. the Reimann Hypothesis). 
\section{Direct Proofs}
We will start with some assumptions, we wlil determine what those assumptions mean, and discern truthness from it.\en
Let's start with some basic assumptions:
\begin{itemize}
    \item An integer $n$ is even if it is possible to write $n=2k$ for some integer $k$.
    \item An integer $n$ is odd if it is possible to write $n = 2\ell + 1$ for some integer $\ell$
\end{itemize}
Quick questions:
\begin{enumerate}
    \item Why is $n= 10$ even? \en
    Because we can write $2\cdot 5 = 10$
    \item Why is $m = -21$ odd? \en
    Because we can write $-11\cdot 2 + 1 = -21$. 
\end{enumerate}
\subsection{First Proof}
Now, consider the following statement:
\centerit{If $n$ is even, then $n^2+1$ is odd.}
Without proof, I think the following statement is true.
\begin{proof}
Assume that $n=2k$. We must show that $n^2 +1$ can be written as $2\ell + 1$. Since $n^2 + 1 = (2k)^2 + 1$, we can write:
\begin{align*}
    4k^2 +1 = 2(2\cdot k^2) + 1
\end{align*}
Let $2\cdot k^2 \equiv \ell$, then we have that:
\begin{align*}
    n^2+1 = 2\ell + 1
\end{align*}
Thus, we see that $n^2+1$ is odd as I so hungrily desired.
\end{proof}
\pagebreak
\subsection{Structure of a proof}
We will quickly overview the general structure of a proof.
\begin{proof}
    Let $n$ be an even integer.\\
    We must show $n^2+1$ is odd. \\
    Since $n$ is even, then $n = 2k, k\in \Z$.\\
    Then $n^2 + 1 = (2k)^2 + 1 = 2(2k^2) + 1$. Let $\ell = 2k^2$, which is an integer. \\
    Then, $n^2 + 1 = 2\ell +1$ and is thus odd, as desired. 
\end{proof}
\subsection{Second Proof}
Fill in the following statement:
\centerit{If $m$ is off, then $m^2 +m $ is even}
\begin{proof}
Assume that $m \equiv 2k + 1$. Then we can write:
\begin{align*}
    m^2+m &= (2k + 1)^2 + 2k + 1\\
    &= 4k^2 + 4k + 1 + 2k +1\\
    &= 4k^2 + 6k + 2.
\end{align*}
Simplifying the equation gives that:
\begin{align*}
    m^2 + m = 2(2k^2 + 3k + 2)
\end{align*}
If we define $\ell \equiv 2k^2 + 3k + 2$, then we write $m^2 + m = 2\ell$. Thus, because we can write $m^2 +m$ in this form, we see that $m^2+m$ is even, as desired. 
\end{proof}
\subsection{Third Proof}
If $n$ is odd, then $n^2 - 1$ is a multiple of 8.
\begin{proof}
    Assume that we may write $n$ as $n = 2k+1, \,k\in \Z$ because $n$ is odd. We must show that $n^2 - 1$ is a multiple of 8, such that we can write $n^2- 1 = 8\ell$, for some $\ell \in \Z$. \en
    Consider that:
    \begin{align*}
        n^2 - 1 &= (2k +1)^2 -1\\
        &= 4k^2 + 4k\\
        &= 4(k^2 + k)
    \end{align*}
    Via proof \#2, we know that $k^2 + k$ is even, so we can write:
    \begin{align*}
        4(k^2+k) &= 4(2\ell)\\
        &= 8\ell, \,\text{ for some } \ell \in \Z
    \end{align*}
    Thus, we have that $n^2 - 1 = 8\ell$, and $n^2 - 1$ is a multiple of 8, as desired.
\end{proof}
\subsection{Fourth Proof}
If $n<10$, then $n^2+100 \le 181$. 
\begin{proof}
    \textit{By counterexample. }We want to show that $\exists n<10, $ for $n\in\Z$ such that $n^2 + 100 > 181$.\en
    Consider $n = -10$. Then we have that:
    \begin{align*}
        n^2 + 100 = 200 > 181
    \end{align*}
    Thus, the statement ``if $n<10$, then $n^2+100 \le 181$'' is false. 
\end{proof}
\newpage
\section{Divisibility \& Vocab}
We say that:
\centerit{\fbox{An integer $b$ is divisible by an integer $a\ne 0$ if $b = a\cdot k$ for some $K\in \Z$}}
The notation we use is $a\mid b$,\;\;(a ``divides'' b).
\subsection{Theorem 2.2}
Fill in the missing parts o the proof:
\begin{proof}
    Assume $a,b,c,d$ are integers and \textbf{nonzero}. Since $a\mid b$ there is an integer $x$ such that $b= x\cdot a$. Similarly, since $c\mid d$, there is an integer $y$ such that $d = y\cdot c$. It follows that $bd = xy\cdot ac$, and therefore $ac\mid bd$. 
\end{proof}
\subsection{Division Algorithm}
Let $n, m \in \Z,\,(m\ne 0)$. \\
Then there exists a unique $q, r\in\Z$ where $n = q\cdot m + r$, for $0\le r < q$
\subsection{Divisibility Proof 1}
Let $n$ be an integer. \textbf{Claim:} it is impossible for $n^2 - 2$ to be divisible by 3. \\
Consider the cases where we can represent $n^2 - 2$ as $3a + r$, where $r = 0,1,2$. Thus we simplify this representation with $n = 0, 1, 2$.  We get $-2, -1, 2$ as our results, all of which are not divisible by 3. 

\end{document}
